\section{Conclusion}
\label{sec:conclusion}

Taken together, the attention and probing results suggest that chemical information is not represented uniformly across mechanisms or layers in these models. Attention exhibits only weak and highly selective alignment with molecular geometry, indicating that spatial structure is not encoded through a small set of geometry-tracking heads. In contrast, probing and causal intervention reveal that chemically meaningful properties are strongly represented within the residual stream, but become computationally usable at substantially different depths across models. More broadly, our results move beyond asking whether chemical information is present, toward understanding how chemically relevant structure is internally organized and used during computation.

\para{Future Directions.}
One natural next step is to study the internal geometry of chemical representations more directly. While linear probing identifies which properties are recoverable, it does not characterize how different chemical features are arranged relative to one another within representation space, whether they occupy shared or independent subspaces, or how these structures evolve across layers. Future work could examine feature overlap, representation superposition, and layerwise organization using methods such as sparse autoencoders, subspace analysis, or representation similarity techniques.

A second direction is to trace how chemical information flows through the network during computation. Our causal interventions show when a representation becomes functionally usable, but do not yet identify which downstream components consume that information or how multiple chemical features are composed into higher-level reasoning. Studying attention circuits, residual-stream interactions, and feature propagation across layers may help reveal how models combine local chemical descriptors into more global molecular representations.

Finally, an important open question is how different pretraining choices shape the resulting internal representations. Although \molm and \smi achieve similar downstream performance, they differ substantially in when and how chemical information becomes usable. Extending this analysis to additional architectures, objectives, and training datasets may help identify which design choices lead to more structured, transferable, or computationally efficient chemical representations.